\documentclass[a4paper, 11pt]{article}
\usepackage{comment} % habilita el uso de comentarios en varias lineas (\ifx \fi) 
\usepackage{lipsum} %Este paquete genera texto del tipo  Lorem Ipsum. 
\usepackage{fullpage} % cambia el margen

\begin{document}
%Construccion del encabezado, asegurate de cambiar tus datos !!!!

oindent
\large\textbf{Actividad 4} \\
\textbf{Analisis de Gestores} \\

ormalsize Investigacion Aplicada a la Ingeniería Ciclo 01-2019 \\
Alumno: ApellidoP ApellidoM, Nombre\hfill Fecha de Reporte: XX/XX/XX \\
 

\section*{Enunciado del Problema}
Coloca aquí el enunciado del problema. Ejemplo de una cita\cite[p.219]{Robotics}. Aquí está otro ejemplo\cite{Flueck}

\section*{Palabras Clave}
\lipsum[2]

\section*{Hipotesis de Solucion}
\lipsum[3]

% Analisis de Gestores Bibliográficos, no mayor a 200 palabras.
\section*{Análisis de Gestores Bibliográficos}
\lipsum[4]
% para comentar secciones, usa el comando \ifx y \fi. Usa esta tecnica cuando escribas tu primer borrador. Por ejemplo, para comentar algo se usa:
%  \ifx
%	\begin{itemize}
%		\item item1
%		\item item2
%	\end{itemize}	
%  \fi

\section*{Conclusiones}
\lipsum[5]



\begin{thebibliography}{9}
\bibitem{Robotics} Fred G. Martin \emph{Robotics Explorations: A Hands-On Introduction to Engineering}. New Jersey: Prentice Hall.
\bibitem{Flueck}  Flueck, Alexander J. 2005. \emph{ECE 100}[online]. Chicago: Illinois Institute of Technology, Electrical and Computer Engineering Department, 2005 [cited 30
August 2005]. Available from World Wide Web: (http://www.ece.iit.edu/~flueck/ece100).
\end{thebibliography}

\end{document}
